\section{The autonomous-holdover floor and the provider channel}
\label{sec:holdover}

The gauge \eqref{eq:gauge} has a temporal part, and that part is what a self-referential
network cannot hold: the common clock offset and, more consequentially, the common clock
\emph{rate}. This section states the holdover floor as a corollary of the temporal gauge,
frames it as the ensemble-time free parameter of the network, and connects the provider
channel of Section~\ref{sec:providersplit} to the external-tie requirement.

\subsection{The temporal gauge and the holdover floor}
\label{sec:holdovercor}

\begin{corollary}[Autonomous-holdover temporal-gauge floor]
\label{cor:holdover}
In the self-referential model, where elapsed time is measured by the network's \emph{own}
drifting clocks, the temporal part of the gauge $\Nsp(\Gmat)$ has dimension two---the common
offset $\dtau$ and the common rate $\dalpha$---so both are internally unobservable
(\Proven{} structure; the self-referential timekeeping setup is \Modelled{}). Adding a single
external rate tie reduces the temporal gauge to dimension one: the common rate becomes
observable and only the common offset survives (\Proven{}, by the PSD intersection law). No
internal measurement removes the common rate; per-node \emph{relative} rates are observable
across the integration windows, but the \emph{common} rate is not.
\end{corollary}

\begin{proof}
Adding a common offset $\dtau$ and a common rate $\dalpha$ to every node's clock changes no
inter-node differential observable, so $\bm e_{\dtau}^{\mathrm{common}}$ and
$\bm e_{\dalpha}^{\mathrm{common}}$ lie in $\Nsp(\Gmat)$ exactly (the temporal part of
\eqref{eq:gauge}), giving temporal-gauge dimension two; the engine reports
\code{rate\_in\_gauge}${}=\text{true}$, \code{temporal\_gauge\_dim}${}=2$. An external rate
tie contributes a PSD information block $\bm{I}_{\mathrm{rt}}$ that annihilates every vector
except those with a common-rate component; by the intersection law $\Nsp(\Gmat^\top\Wmat\Gmat
+\bm{I}_{\mathrm{rt}})=\Nsp(\Gmat^\top\Wmat\Gmat)\cap\Nsp(\bm{I}_{\mathrm{rt}})$, the common
rate is removed while the common offset survives, so \code{rate\_in\_gauge}${}=\text{false}$,
\code{temporal\_gauge\_dim}${}=1$. Multi-window observations constrain \emph{differences} of
node rates (per-node relative rates) but not their common part, which cancels differentially.
\end{proof}

The magnitude of this floor is the free-running physical drift of the network's clock
ensemble relative to Lunar Coordinate Time (TCL)---a \Modelled{} quantity set by the
oscillators, not by geometry. What is \Proven{} is the \emph{structure}: the common rate is a
genuine gauge direction with zero internal information (not merely a weakly-observed
direction), so an autonomous monitor cannot detect or correct it. This sharpens, in the
self-referential setting, the ensemble-time free-parameter idea that a timescale realized by
an ensemble of clocks has an irreducible common drift only an external reference fixes.

\subsection{The ensemble-time free parameter and Earth-link-loss holdover}
\label{sec:ensemble}

The holdover floor is precisely the contingency the paper is framed around. Under nominal
operation a lunar service is Earth-referenced---an Earth--Moon time comparison is exactly the
external rate tie of Corollary~\ref{cor:holdover}, pinning the common rate. When the Earth
link is lost, that tie is gone: the network keeps time with its own clocks, and the common
rate re-enters the gauge as an unobservable, uncorrectable drift. This is the temporal
counterpart of the datum-defect story---a lunar coordinate time realized internally has a
free rate parameter, just as a lunar frame realized internally has free rigid directions
\citep{baweja2026gauge}. Two current strands of lunar time standardization frame this floor,
and we keep them distinct. \emph{Lunar Coordinate Time} (TCL) is the IAU 2024 Resolution~II
\emph{coordinate} time---a TCG-analog defined for a lunar celestial reference system
\citep{iau2024ltc}---and is the ideal against which the network's internal drift is measured;
\emph{Coordinated Lunar Time} (LTC) is the \emph{realized} standard, a UTC-analog motivated by
the April 2024 US OSTP memorandum \citep{iau2024ltc,fienga2024moonlight}, that an operational
lunar service must actually disseminate. Either way the floor matters: a realized lunar
timescale must specify the external tie and cadence that bound the holdover, because no
internal monitor can. We state this as \emph{relevance}, not endorsement: nothing here is a
certified timescale realization or an agency-approved holdover budget.

\subsection{The provider channel}
\label{sec:providerchannel}

The provider fault is the structured, multi-node channel by which an autonomous network can
still cross-check its own frame and timescale without an Earth link---but only partly.
Proposition~\ref{prop:provider} splits a provider-wide bias into a common part in
$\rng(\Gmat)$ and a differential part outside it. The common part is exactly the class an
external tie is needed for: it is the shared component of the inter-ephemeris orientation and
position floor of P2 \citep{baweja2026interop}, and, being a whole-network datum move, it is
invisible to any number of mutually-consistent providers. The differential part is what
autonomy \emph{can} exploit: two providers detect their disagreement (but cannot say which is
wrong---the two-clock problem), and three or more identify the culprit through the block-spark
bound of Theorem~\ref{thm:byzantine}. The provider-count trend of Table~\ref{tab:provider}
($1\to0$, $2\to48$, $4\to72$ detectable cross-provider rows) is the \Modelled{} quantitative
face of this split. The design consequence---engineer at least three independent providers to
identify a faulty one, and retain an external tie for the common part---is
Section~\ref{sec:designlaws}.
